% GENERATED FILE. Edit canonical lessons, then run npm run generate:books.
% canonical-source-hash: fnv1a64:3e5769a530e24ceb
% canonical-lessons: IT-C26-primo, IT-C26-secondo, IT-C26-terzo, IT-C26-quarto, IT-C26-quinto

\chapter{Primo, Secondo, Terzo}
\label{ch:it-ordinali}
\hlchaptermodality{26}

\begin{chapteropening}
\textbf{By the end of this chapter:} \emph{I can say first to fifth in Italian with the right ending, read the two courses on an Italian menu, and ask for a quarter of an hour.}

The fifth ordinal said against the four before it, and the sound split that makes cinque and quinto refuse to rhyme.
\end{chapteropening}

\section[primo / prima]{primo — first, and a word this book already used}

\label{lesson:IT-C26-primo}



\begin{quote}
A new kind of counting starts here, so the recalls come first, at the four growing distances this book measures.

\begin{itemize}
\raggedright
  \item \emph{Recall:} say \emph{la gola} — \textbf{R1}, one lesson back, before it decays
  \item \emph{Recall:} say \emph{il caffè} — \textbf{R2}, the first real retrieval, five to fifteen lessons back
  \item \emph{Recall:} say \emph{abitare} — \textbf{R3}, durable, twenty to sixty lessons back
  \item \emph{Recall:} say \emph{uno, due, tre, quattro, cinque} — \textbf{R3}, a second word from the same distance
  \item \emph{Recall:} say \emph{ciao} — \textbf{R4}, recognition at distance, eighty lessons back or more
  \item \emph{Recall:} say \emph{buono / buon} — \textbf{R4}, a second word from the same distance
\end{itemize}
\end{quote}

\subsection*{You'll want to know: The word}
\begin{quote}
\textbf{primo} (masculine) · \textbf{prima} (feminine) — \textbf{first}. \textbf{il primo giorno} — the first day. \textbf{la prima volta} — the first time.
\end{quote}

\emph{Uno, due, tre} answer \textbf{how many}. This word answers \textbf{which one in the line}, and Italian keeps a second series for that job. It behaves as any describing word does and changes its ending with what it counts.

\begin{cousinweb}
This book has used \emph{primo} once, and hidden it. The seasons lesson glossed \textbf{primavera} as \textquotedblleft{}the \textbf{first} green\textquotedblright{} — \emph{prima vera}, the first spring — and the word behind that gloss was never given. Now it is.

One thing to know before you use it. \textbf{Il primo piano} is not the ground floor. Italian calls the ground floor \textbf{il pianterreno}, \textquotedblleft{}the earth floor\textquotedblright{}, and starts counting at the top of the first flight of stairs — so \emph{il primo piano} is what an American calls the second floor. A lift panel in Rome begins \textbf{T}, then \textbf{1}.
\end{cousinweb}

\subsection*{Your turn}
Say these aloud:

\begin{itemize}
\raggedright
  \item \emph{primo, prima}
  \item \emph{il primo giorno} — the first day
  \item \emph{primavera} — the first spring
\end{itemize}

\begin{tcolorbox}[breakable,colback=teal!4,colframe=teal!35!black,title={Before you move on}]
Say \textquotedblleft{}the first time.\textquotedblright{} (\textbf{La prima volta}.) Which word in this book was already built on \emph{primo}? (\emph{\textbf{Primavera}} — the first spring.) And which floor is \emph{il primo piano}? (\textbf{The one above the ground floor}, which Italian calls \emph{il pianterreno}.)
\end{tcolorbox}
\section[secondo / seconda]{secondo — the one that follows, and half a menu}

\label{lesson:IT-C26-secondo}



\begin{quote}
Recalls before the second ordinal, at the four measured distances.

\begin{itemize}
\raggedright
  \item \emph{Recall:} say \emph{primo} — \textbf{R1}, one lesson back, before it decays
  \item \emph{Recall:} say \emph{il tè} — \textbf{R2}, the first real retrieval, five to fifteen lessons back
  \item \emph{Recall:} say \emph{lavorare} — \textbf{R3}, durable, twenty to sixty lessons back
  \item \emph{Recall:} say \emph{sei, sette, otto, nove, dieci} — \textbf{R3}, a second word from the same distance
  \item \emph{Recall:} say \emph{il / la / lo} — \textbf{R4}, recognition at distance, eighty lessons back or more
  \item \emph{Recall:} say \emph{giorno} — \textbf{R4}, a second word from the same distance
\end{itemize}
\end{quote}

\subsection*{You'll want to know: The word}
\begin{quote}
\textbf{secondo} (masculine) · \textbf{seconda} (feminine) — \textbf{second}. \textbf{il secondo giorno} — the second day. \textbf{la seconda volta} — the second time.
\end{quote}

There is no \emph{due} in it. It comes from Latin \emph{\textbf{sequī}}, \textquotedblleft{}\textbf{to follow}\textquotedblright{}, so second place is named for what it does — the place that \textbf{comes after} — and not for the count it sits at. English \textbf{second}, \textbf{sequel}, \textbf{sequence} and \textbf{consequence} are the same root.

\begin{culture}
An Italian menu is written with these two words used as nouns, and once you have them you can read one:

\noindent
\begin{tabularx}{\linewidth}{@{}>{\raggedright\arraybackslash}X>{\raggedright\arraybackslash}X@{}}
\toprule
\textbf{course} & \textbf{what arrives} \\
\midrule
\textbf{l'antipasto} & before the meal \\
\textbf{il primo} & pasta, rice or soup \\
\textbf{il secondo} & meat or fish \\
\textbf{il dolce} & the sweet \\
\bottomrule
\end{tabularx}

\emph{Il primo} and \emph{il secondo} are not \textquotedblleft{}starter\textquotedblright{} and \textquotedblleft{}main\textquotedblright{}. They are literally \textbf{the first} and \textbf{the second}, and an Italian ordering both is ordering two courses, not a starter and an entrée. Two ordinals, and half a meal is legible.
\end{culture}

\subsection*{Your turn}
Say these aloud:

\begin{itemize}
\raggedright
  \item \emph{secondo, seconda}
  \item \emph{la seconda volta} — the second time
  \item the menu — \emph{antipasto, primo, secondo, dolce}
\end{itemize}

\begin{tcolorbox}[breakable,colback=teal!4,colframe=teal!35!black,title={Before you move on}]
What verb is \emph{secondo} built on? (\emph{\textbf{Sequī}}, to follow.) What are \emph{il primo} and \emph{il secondo} on a menu? (\textbf{The first and second courses} — pasta, then meat.) And say \textquotedblleft{}the second day.\textquotedblright{} (\textbf{Il secondo giorno}.)
\end{tcolorbox}
\section[terzo / terza]{terzo — where the number starts to show through}

\label{lesson:IT-C26-terzo}



\begin{quote}
Recalls first, then the third ordinal — the first with a number visible inside it.

\begin{itemize}
\raggedright
  \item \emph{Recall:} say \emph{secondo} — \textbf{R1}, one lesson back, before it decays
  \item \emph{Recall:} say \emph{il latte} — \textbf{R2}, the first real retrieval, five to fifteen lessons back
  \item \emph{Recall:} say \emph{the fifth chapter's own recombining} — \textbf{R3}, durable, twenty to sixty lessons back
  \item \emph{Recall:} say \emph{lunedì, martedì, mercoledì} — \textbf{R3}, a second word from the same distance
  \item \emph{Recall:} say \emph{buongiorno} — \textbf{R4}, recognition at distance, eighty lessons back or more
  \item \emph{Recall:} say \emph{buonasera} — \textbf{R4}, a second word from the same distance
\end{itemize}
\end{quote}

\subsection*{You'll want to know: The word}
\begin{quote}
\textbf{terzo} (masculine) · \textbf{terza} (feminine) — \textbf{third}. \textbf{il terzo piano} — the third floor. \textbf{la terza volta} — the third time.
\end{quote}

Say \emph{tre} and then \emph{terzo}. Two ordinals have gone past with no number in them at all; this is the first where you can hear one. The \emph{z} is \textbf{ts} — \emph{TER-tso}.

\begin{cousinweb}
\textbf{Un terzo} also means \textquotedblleft{}\textbf{a third}\textquotedblright{}, the fraction, and that is the same word used as a noun. Italian names its small fractions with the ordinals:

\noindent
\begin{tabularx}{\linewidth}{@{}>{\raggedright\arraybackslash}X>{\raggedright\arraybackslash}X@{}}
\toprule
\textbf{fraction} & \textbf{Italian} \\
\midrule
a third & \textbf{un terzo} \\
a quarter & \textbf{un quarto} \\
a fifth & \textbf{un quinto} \\
\bottomrule
\end{tabularx}

So the words this chapter is teaching do a second job you will meet in a recipe or on a bill. English does exactly the same thing with \emph{a third} and \emph{a quarter}, and for the same reason: the ordinal already names a position in a division.
\end{cousinweb}

\subsection*{Your turn}
Say these aloud:

\begin{itemize}
\raggedright
  \item \emph{terzo, terza} — \emph{TER-tso}
  \item the join — \emph{tre $\to$ terzo}
  \item \emph{un terzo} — a third
\end{itemize}

\begin{tcolorbox}[breakable,colback=teal!4,colframe=teal!35!black,title={Before you move on}]
Which number is inside \emph{terzo}? (\emph{\textbf{Tre}}.) How is the \emph{z} said? (\textbf{Ts} — \emph{TER-tso}.) And what does \emph{un terzo} mean besides \textquotedblleft{}the third one\textquotedblright{}? (\textbf{A third} — the fraction.)
\end{tcolorbox}
\section[quarto / quarta]{quarto — the fourth, and a quarter of an hour}

\label{lesson:IT-C26-quarto}



\begin{quote}
Recalls first, then the fourth — and the phrase that makes these words immediately useful.

\begin{itemize}
\raggedright
  \item \emph{Recall:} say \emph{terzo} — \textbf{R1}, one lesson back, before it decays
  \item \emph{Recall:} say \emph{lo zucchero} — \textbf{R2}, the first real retrieval, five to fifteen lessons back
  \item \emph{Recall:} say \emph{sabato e domenica} — \textbf{R3}, durable, twenty to sixty lessons back
  \item \emph{Recall:} say \emph{l'ora} — \textbf{R3}, a second word from the same distance
  \item \emph{Recall:} say \emph{buonanotte} — \textbf{R4}, recognition at distance, eighty lessons back or more
  \item \emph{Recall:} say \emph{grazie} — \textbf{R4}, a second word from the same distance
\end{itemize}
\end{quote}

\subsection*{You'll want to know: The word}
\begin{quote}
\textbf{quarto} (masculine) · \textbf{quarta} (feminine) — \textbf{fourth}. \textbf{il quarto piano} — the fourth floor. \textbf{la quarta volta} — the fourth time.
\end{quote}

\emph{Quattro} is four and \emph{quarto} is four with the ordinal ending, the double \emph{t} squeezed out. Two in a row now behave the same way.

\subsection*{You'll want to know: a quarter of an hour}
\begin{quote}
\textbf{un quarto d'ora} — a quarter of an hour, fifteen minutes. \textbf{tre quarti d'ora} — three quarters of an hour.
\end{quote}

The clock chapter gave you \textbf{l'ora}. This is what Italian does with it, and it is the single most useful thing in the chapter: \emph{fra un quarto d'ora}, \textquotedblleft{}in a quarter of an hour\textquotedblright{}, is how somebody tells you to wait a bit.

It also gives you \textbf{le sei e un quarto} — a quarter past six — where the ordinal is doing the work an English speaker expects from \textquotedblleft{}quarter\textquotedblright{}. Same idea, and the word you need for it is one you now have.

\subsection*{Your turn}
Say these aloud:

\begin{itemize}
\raggedright
  \item \emph{quarto, quarta}
  \item \emph{un quarto d'ora} — a quarter of an hour
  \item \emph{le sei e un quarto} — a quarter past six
\end{itemize}

\begin{tcolorbox}[breakable,colback=teal!4,colframe=teal!35!black,title={Before you move on}]
Say \textquotedblleft{}a quarter of an hour.\textquotedblright{} (\textbf{Un quarto d'ora}.) What number is inside \emph{quarto}? (\emph{\textbf{Quattro}}.) And how would you say \textquotedblleft{}a quarter past six\textquotedblright{}? (\textbf{Le sei e un quarto}.)
\end{tcolorbox}
\section[quinto / quinta]{quinto — the fifth, and a sound the cardinal lost}

\label{lesson:IT-C26-quinto}



\begin{quote}
Recalls first, then the fifth — and a pair that does not rhyme when it ought to.

\begin{itemize}
\raggedright
  \item \emph{Recall:} say \emph{quarto} — \textbf{R1}, one lesson back, before it decays
  \item \emph{Recall:} say \emph{l'amico e l'amica} — \textbf{R2}, the first real retrieval, five to fifteen lessons back
  \item \emph{Recall:} say \emph{la gola} — \textbf{R2}, a second word from the same distance
  \item \emph{Recall:} say \emph{mezzogiorno e mezzanotte} — \textbf{R3}, durable, twenty to sixty lessons back
  \item \emph{Recall:} say \emph{i mesi} — \textbf{R3}, a second word from the same distance
  \item \emph{Recall:} say \emph{prego} — \textbf{R4}, recognition at distance, eighty lessons back or more
  \item \emph{Recall:} say \emph{prego, go ahead} — \textbf{R4}, a second word from the same distance
\end{itemize}
\end{quote}

\subsection*{You'll want to know: The word}
\begin{quote}
\textbf{quinto} (masculine) · \textbf{quinta} (feminine) — \textbf{fifth}. \textbf{il quinto piano} — the fifth floor. \textbf{la quinta volta} — the fifth time.
\end{quote}

\begin{sounds}
Say them together: \emph{\textbf{cinque}}, \emph{\textbf{quinto}}. They are the same Latin root and they do not sound alike, and the difference is worth a minute.

Latin had \emph{quīnque} and \emph{quīntus}, both beginning \textbf{kw-}. Then the first one was pulled forward by the \emph{i} that followed and softened all the way to \textbf{ch-} — which is why \emph{cinque} opens like English \emph{chin}. The ordinal was not pulled, and it still opens \textbf{kw-}: \emph{KWIN-to}.

So this pair is a sound change caught halfway through: one word moved and its twin did not. Listen for the same split in \emph{dieci} against \emph{decimo}, five lessons ahead — the cardinal softened, the ordinal did not.
\end{sounds}

\subsection*{Your turn}
Say these aloud:

\begin{itemize}
\raggedright
  \item \emph{quinto, quinta} — \emph{KWIN-to}
  \item the pair — \emph{cinque} soft, \emph{quinto} hard
  \item \emph{la quinta volta} — the fifth time
\end{itemize}

\begin{tcolorbox}[breakable,colback=teal!4,colframe=teal!35!black,title={Before you move on}]
Why do \emph{cinque} and \emph{quinto} not rhyme? (\textbf{The cardinal softened its opening to \emph{ch-} and the ordinal kept \emph{kw-}.}) Say the fifth floor. (\textbf{Il quinto piano}.) And which pair five lessons ahead does the same thing? (\emph{\textbf{Dieci}} against \emph{\textbf{decimo}}.)
\end{tcolorbox}
